% Apresentação: Evolução do Portfólio de Políticas (Julho -> Agosto/2026)
% Antes/depois de cada política, defeitos corrigidos e novas famílias.
% Compilar com: latexmk -xelatex main.tex  (ver Makefile)
\documentclass[11pt,aspectratio=169]{beamer}

\usetheme{metropolis}

% --- Idioma e fontes ---
\usepackage{fontspec}
\usepackage{polyglossia}
\setdefaultlanguage[variant=brazilian]{portuguese}

% --- Pacotes de suporte ---
\usepackage{booktabs}
\usepackage{makecell}
\usepackage{colortbl}
\usepackage{array}
\usepackage{graphicx}
\usepackage{tikz}
\usetikzlibrary{positioning,arrows.meta,fit,backgrounds,calc}
\usepackage{amsmath}
\usepackage{textcomp}

% --- Paleta acadêmica (idêntica à apresentação de qualificação) ---
\definecolor{darkblue}{HTML}{0F2F57}
\definecolor{slategray}{HTML}{4D6175}
\definecolor{lightgray}{HTML}{EAF0F6}
\definecolor{accenteal}{HTML}{0E8C95}
\definecolor{accentpetrol}{HTML}{1E5D8C}
\definecolor{warnred}{HTML}{9A3A3A}
\definecolor{softink}{HTML}{1F3348}
\definecolor{cardbg}{HTML}{F4F8FC}
\definecolor{canvasbg}{HTML}{FFFFFF}
% --- Cores extras: antes (defeito/legado) vs depois (corrigido/novo) ---
\definecolor{beforered}{HTML}{9A3A3A}
\definecolor{aftergreen}{HTML}{1E7A4C}

\setbeamercolor{background canvas}{bg=canvasbg}
\setbeamercolor{normal text}{fg=softink,bg=white}

\setbeamercolor{palette primary}{bg=darkblue,fg=white}
\setbeamercolor{palette secondary}{bg=accenteal!85!darkblue,fg=white}
\setbeamercolor{palette tertiary}{bg=darkblue,fg=white}
\setbeamercolor{frametitle}{bg=white,fg=darkblue}
\setbeamercolor{progress bar}{fg=accenteal,bg=darkblue!12}
\setbeamercolor{alerted text}{fg=accentpetrol!85!darkblue}
\setbeamercolor{item}{fg=accenteal!90!darkblue}
\setbeamercolor{subitem}{fg=accentpetrol!85!darkblue}

\setbeamerfont{title}{size=\Large,series=\bfseries}
\setbeamerfont{subtitle}{size=\normalsize}
\setbeamerfont{frametitle}{size=\large,series=\bfseries}
\setbeamerfont{block title}{series=\bfseries}
\setbeamerfont{normal text}{size=\normalsize}

\setbeamertemplate{frame numbering}[fraction]
\metroset{block=fill,sectionpage=none,numbering=fraction,progressbar=frametitle}

\setbeamercolor{block title}{fg=darkblue,bg=lightgray}
\setbeamercolor{block body}{fg=softink,bg=cardbg}
\setbeamercolor{alertblock title}{fg=white,bg=accentpetrol!65!darkblue}
\setbeamercolor{alertblock body}{fg=softink,bg=accentpetrol!8}

\setbeamertemplate{headline}{
  \begin{beamercolorbox}[wd=\paperwidth,ht=0.14em,dp=0ex]{progress bar}
  \end{beamercolorbox}
}

\setbeamertemplate{section page}{
  \begin{centering}
    \vspace*{0.17\paperheight}
    {\color{darkblue}\rule{0.42\linewidth}{1.1pt}}\par
    \vspace{0.9em}
    {\usebeamerfont{section title}\color{darkblue}\insertsectionhead\par}
    \vspace{0.9em}
    {\color{accentpetrol}\rule{0.24\linewidth}{0.9pt}}\par
    \vfill
  \end{centering}
}

\AtBeginSection[]{\begin{frame}[plain,noframenumbering]\sectionpage\end{frame}}

% --- Comandos auxiliares de estilo ---
\newcommand{\highlight}[1]{\textcolor{accenteal}{\textbf{#1}}}
\newcommand{\fonte}[1]{{\footnotesize\textcolor{slategray}{Fonte: #1}}}
\newcommand{\keymessage}[1]{%
  \begin{beamercolorbox}[wd=\linewidth,rounded=true,sep=0.8em,leftskip=0.8em,rightskip=0.8em]{block body}
    \textbf{\textcolor{darkblue}{#1}}
  \end{beamercolorbox}%
}
\newcommand{\tabstyle}{\renewcommand{\arraystretch}{1.22}}

% --- Comandos específicos deste deck: rótulos ANTES (julho) / DEPOIS (atual) ---
\newcommand{\antes}{\textcolor{beforered}{\textbf{Antes (jul/2026)}}}
\newcommand{\depois}{\textcolor{aftergreen}{\textbf{Depois (atual)}}}
\newcommand{\novo}{\textcolor{aftergreen}{\textbf{Novo (ago/2026)}}}

\title{Da Implementação de Julho ao Portfólio Atual}
\subtitle{O que mudou em cada uma das 18 políticas de reposição do AIPE}
\author{Matheus Inácio Batista Santos}
\institute{\small Universidade Federal de Alagoas (UFAL)\\
\small Instituto de Computação (IC)\\
\scriptsize Programa de Pós-Graduação em Informática (PPGI)\\
\scriptsize Orientadores: Dr. Rian G. S. Pinheiro e Dr. Bruno C. S. Nogueira}
\date{\small Apresentação técnica de acompanhamento}

\begin{document}

% Slide 1 — Título
\begin{frame}[plain]
  \titlepage
  \note{
    Esta apresentação documenta a evolução do portfólio de políticas de
    reposição do AIPE entre a versão usada no Experimento 1 (julho/2026) e a
    reimplementação concluída em agosto/2026. O objetivo é mostrar,
    política por política, exatamente o que mudou: dois defeitos
    corrigidos (PPO e escala da aptidão das meta-heurísticas), seis
    políticas novas, vetorização do simulador e a extensão da função de
    aptidão com termos de risco. Nenhum resultado numérico novo é
    reportado aqui -- o foco é a implementação.
  }
\end{frame}


\begin{frame}{Sumário}
  \begin{enumerate}
    \item Por que reimplementar?
    \item O simulador: de laço Python a lote vetorizado
    \item Políticas clássicas: o que ficou e o que entrou
    \item Meta-heurísticas: GA, SA, PSO, DE
    \item A função de aptidão: de soma ponderada a restrição com risco
    \item Aprendizado por reforço: DQN, PPO, SARSA
    \item Híbridas GA-RL
    \item Infraestrutura e hiperparâmetros
    \item Portfólio final e próximos passos
  \end{enumerate}
\end{frame}

\section{Por que reimplementar?}
% Slide — Por que reimplementar
\begin{frame}{Por que reimplementar o portfólio inteiro?}
  \textbf{\textcolor{darkblue}{Não foi troca de biblioteca. Foram dois defeitos que afetam resultados já reportados.}}
  \vspace{0.4em}

  \begin{itemize}
    \item A política rotulada \textbf{PPO não otimizava o objetivo do PPO}: a
      atualização do ator estava matematicamente errada (próximo slide).
    \item A \textbf{aptidão das meta-heurísticas} (GA/SA/PSO/DE) usava soma
      ponderada com pesos fixos -- ponto de operação sensível à escala do
      custo, que varia por ordens de grandeza entre séries.
  \end{itemize}
  \vspace{0.6em}

  \textbf{Consequência prática:} qualquer correção pontual exigiria retreinar
  RL e recalibrar meta-heurísticas de qualquer forma. A decisão foi
  aproveitar a reescrita para também:
  \begin{itemize}
    \item vetorizar o simulador (laço Python $\rightarrow$ lote em PyTorch);
    \item ampliar o portfólio de 12 para \highlight{18 políticas}, cobrindo
      estado da arte publicado onde ele existe;
    \item alinhar hiperparâmetros a uma referência auditável (Zabraoui et
      al., 2025) onde não existe estado da arte de meta-heurísticas para
      inventário.
  \end{itemize}

  \note{
    Esta reimplementação não nasceu de uma preferência por PyTorch. Nasceu
    de dois defeitos concretos encontrados ao auditar o código depois do
    Experimento 1: o PPO tinha o gradiente do ator errado, e a aptidão das
    meta-heurísticas misturava níveis de serviço e custo numa soma
    ponderada cuja escala dependia do volume da série. Como corrigir os
    dois exigia mexer no núcleo de RL e de meta-heurísticas de qualquer
    forma, a decisão foi reescrever o simulador de forma vetorizada ao
    mesmo tempo -- o que abriu espaço para rodar populações inteiras em
    uma única simulação, tornando viável ampliar o portfólio sem explodir
    o tempo de execução.
  }
\end{frame}

% Slide — Os dois defeitos, em detalhe
\begin{frame}{Os dois defeitos que motivaram a reescrita}
  \tabstyle\footnotesize
  \begin{tabular}{@{}p{0.16\linewidth}p{0.40\linewidth}p{0.40\linewidth}@{}}
    \toprule
    & \antes & \depois \\
    \midrule
    \textbf{PPO} &
      Ator atualizado somando a vantagem $A_t$ diretamente aos logits-alvo e
      treinado por \textbf{erro quadrático}. Isso não é o gradiente da
      \textit{surrogate function} recortada do PPO. &
      Gradiente correto por autograd sobre a razão de verossimilhanças
      $r_t=\pi_{\text{novo}}/\pi_{\text{antigo}}$, com \textbf{clipping}
      explícito e \textbf{GAE}$(\lambda)$. \\
    \addlinespace
    \textbf{Aptidão GA/SA/PSO/DE} &
      Soma ponderada fixa $\lambda_1\,\mathrm{NS}-\lambda_2\,\mathrm{CTI}$.
      O ponto de operação depende da \textbf{escala do CTI}, que varia até
      124$\times$ entre a maior e a menor loja do recorte piloto. &
      Formulação \textbf{restrita} (Eq.~4.2 da proposta):
      $\min \mathrm{CTI}$ s.a. $\mathrm{NS}\geq\alpha_{\min}$, com
      penalidade \textbf{relativa ao próprio CTI} -- adimensional, comparável
      entre séries. \\
    \bottomrule
  \end{tabular}
  \vspace{0.8em}

  \begin{alertblock}{Por que isso importa para os resultados já reportados}
    \small A degeneração observada no Experimento 1 (frequência de pedidos
    FP $=0{,}98$ para a política PPO -- pedir quase todo ciclo) é
    \highlight{consistente com o defeito do gradiente}, não com uma
    limitação real do método PPO.
  \end{alertblock}

  \note{
    O defeito do PPO é sutil de detectar olhando só para o código, mas
    óbvio olhando para o resultado: um agente PPO bem treinado não deveria
    pedir em 98\% dos ciclos -- isso é comportamento de política quase
    aleatória, não de política otimizada. A causa raiz é que a atualização
    do ator não seguia o gradiente do objetivo clipped surrogate do PPO; ela
    empurrava os logits na direção da vantagem e depois ajustava por MSE,
    o que é uma heurística sem garantia de melhoria monotônica. O segundo
    defeito, da aptidão das meta-heurísticas, é mais estrutural: como o
    custo varia por ordens de grandeza entre séries de volumes diferentes,
    o mesmo par de pesos (lambda1, lambda2) que funciona bem numa loja de
    alto giro pode privilegiar demais o serviço numa loja pequena, ou
    demais o custo numa loja grande -- contaminando qualquer comparação
    entre políticas feita com a mesma configuração de pesos.
  }
\end{frame}

% Slide — Escopo da mudança: panorama geral
\begin{frame}{Panorama: o que mudou em cada família}
  \centering
  \resizebox{\linewidth}{!}{%
  \tabstyle
  \begin{tabular}{@{}lll@{}}
    \toprule
    \textbf{Família} & \textbf{O que mudou} & \textbf{Natureza} \\
    \midrule
    Clássicas (EOQ, $(s,S)$, Jornaleiro) & Nada -- mantidas como baseline & inalterada \\
    SOTA clássicas (PIL, CappedBaseStock, BDN) & Não existiam & \textcolor{aftergreen}{nova} \\
    Zabraoui (Min-Max, Fixed Interval, Vendor-Resp.) & Não existiam & \textcolor{aftergreen}{nova} \\
    GA & Vetorizado; +elitismo; aptidão corrigida; \textbf{crossover/mutação
      corrigidos p/ bater c/ o artigo} & \textcolor{beforered}{corrigida + vetorizada} \\
    SA & \texttt{scipy.dual\_annealing} $\to$ Metropolis explícito, cadeias paralelas & \textcolor{accentpetrol}{reescrita} \\
    PSO & Sem constrição $\to$ fator de constrição de Clerc \& Kennedy & \textcolor{beforered}{corrigida} \\
    DE & \texttt{scipy.differential\_evolution} $\to$ DE/rand/1/bin vetorizado & \textcolor{accentpetrol}{reescrita} \\
    DQN & Rede NumPy manual $\to$ Double DQN dueling; \textbf{epsilon/episódios
      corrigidos p/ bater c/ o artigo} & \textcolor{beforered}{reescrita + corrigida} \\
    PPO & \textbf{Gradiente do ator corrigido} + GAE; episódios corrigidos & \textcolor{beforered}{corrigida} \\
    SARSA & Tabular 1D $\to$ Expected SARSA, discretização 2D & \textcolor{accentpetrol}{reescrita} \\
    GA-DQN / GA-PPO & Piso de quantidade (\texttt{floor\_ratio}) adicionado & \textcolor{accentpetrol}{ajustada} \\
    \bottomrule
  \end{tabular}}
  \vspace{0.8em}
  \footnotesize\textcolor{slategray}{Vermelho = corrige defeito. Verde = nova política. Azul-petróleo = reescrita sem mudança de defeito conhecido (vetorização/robustez).}

  \note{
    Este slide é o mapa da apresentação: cada linha vira uma ou mais seções
    adiante, com o antes e o depois em detalhe. Vale notar a assimetria --
    as três políticas clássicas originais (EOQ, s-S, Jornaleiro) não
    mudaram nada, porque continuam sendo o baseline de referência do
    Capítulo 5. Tudo que é novo entra como alternativa moderna do mesmo
    slot conceitual, não como substituição.
  }
\end{frame}


\section{O simulador}
% Slide — Simulador: antes/depois
\begin{frame}{O simulador: de laço Python a lote vetorizado}
  \begin{columns}[T,onlytextwidth]
    \column{0.48\textwidth}
      \antes\\[0.3em]
      \small \texttt{InventoryEnv}: uma trajetória por chamada.
      \begin{itemize}\footnotesize
        \item Um indivíduo/partícula = 1 simulação de $T$ ciclos, em laço Python.
        \item População 100 $\times$ 50 gerações do GA $=$ \textbf{5.000
          simulações sequenciais}.
        \item Cada `step` é uma chamada de função Python por ciclo.
      \end{itemize}
    \column{0.48\textwidth}
      \depois\\[0.3em]
      \small \texttt{BatchInventoryEnv}: $B$ trajetórias em paralelo.
      \begin{itemize}\footnotesize
        \item Toda a população avança \textbf{um ciclo por vez, em bloco}
          (operações tensoriais de largura $B$).
        \item Mesma configuração do GA $=$ \textbf{50 chamadas} de simulação
          de largura 100, não 5.000 chamadas de largura 1.
        \item Custo passa de $O(B\times T)$ chamadas Python para $O(T)$
          operações tensoriais.
      \end{itemize}
  \end{columns}
  \vspace{0.8em}

  \begin{block}{O que NÃO mudou}
    \small Ordem das operações por ciclo (chegada $\to$ demanda $\to$
    custo), \textit{lost sales} (demanda não atendida é perdida, não
    acumulada), pipeline de $L{+}1$ posições, e os cinco KPIs (CTI, NS,
    TR, BE, FP). O simulador vetorizado é uma reimplementação, não uma
    reformulação do ambiente.
  \end{block}

  \note{
    A mudança central não é "trocar numpy por PyTorch". É que a população
    inteira de uma meta-heurística passou a ser avaliada em UMA chamada ao
    simulador, em vez de um laço Python por indivíduo. Isso é o que torna
    GPU potencialmente útil neste projeto -- embora, como veremos adiante,
    a medição real mostrou que não compensa nesta escala de população.
    O ambiente continua representando exatamente as mesmas regras de
    negócio: mesma ordem de operações, mesma política de venda perdida,
    mesmo pipeline de lead time.
  }
\end{frame}

% Slide — Garantia de fidelidade do simulador vetorizado
\begin{frame}{Garantia: o novo simulador reproduz o antigo}
  \keymessage{Reescrever o simulador é o maior risco desta mudança -- por
  isso ele foi validado numericamente antes de qualquer política ser
  reimplementada sobre ele.}
  \vspace{0.8em}

  \begin{itemize}
    \item Teste de equivalência (\texttt{tests/test\_env\_equivalence.py}):
      mesma política, mesma demanda, mesma semente, comparando
      \texttt{InventoryEnv} escalar contra \texttt{BatchInventoryEnv} em
      \textbf{30 casos}.
    \item Erro relativo máximo: $\mathbf{3{,}7\times10^{-7}}$ em CTI, NS,
      TR, BE e FP -- compatível com precisão de ponto flutuante, não com
      diferença de lógica.
  \end{itemize}
  \vspace{0.6em}

  \begin{alertblock}{Consequência metodológica}
    \small Qualquer diferença nos resultados entre a versão de julho e a
    atual é atribuível às políticas e à função de aptidão -- \highlight{não
    a uma mudança silenciosa nas regras do ambiente de simulação}.
  \end{alertblock}

  \note{
    Este é o slide de credibilidade da reescrita. Antes de reimplementar
    qualquer política, o simulador vetorizado foi comparado contra o
    original em 30 casos com política, demanda e semente idênticas -- o
    erro relativo máximo encontrado, 3,7e-7, está na faixa de erro de
    arredondamento de ponto flutuante, não de divergência de lógica. Isso
    significa que qualquer diferença de resultado que apareça mais adiante
    (por exemplo, no comportamento do GA ou do PPO) pode ser atribuída com
    confiança às mudanças na política ou na função de aptidão, e não a uma
    mudança silenciosa nas regras do jogo.
  }
\end{frame}


\section{Políticas clássicas: o que ficou e o que entrou}
% Slide — Clássicas: mantidas intactas
\begin{frame}{As três clássicas originais: nada mudou}
  \centering
  \tabstyle\footnotesize
  \begin{tabular}{@{}lll@{}}
    \toprule
    \textbf{Política} & \textbf{Regra} & \textbf{Parâmetros} \\
    \midrule
    EOQ & $Q^\ast=\sqrt{2D K/h}$; pede quando $I{+}\text{pipeline}\leq \mathrm{ROP}$ &
      $\mathrm{ROP}=\mu L + z\sigma\sqrt{L}$ \\
    $(s,S)$ & Banda min-max: pede até $S$ quando posição $\leq s$ &
      $s=\mu L{+}z\sigma\sqrt{L}$; $S=s{+}Q_{\mathrm{EOQ}}$ \\
    Jornaleiro & $Q^\ast=\hat\mu+z\hat\sigma$ (quantil crítico fixo) &
      $z=1{,}28$ (SL alvo 90\%) \\
    \bottomrule
  \end{tabular}
  \vspace{1em}

  \begin{block}{Por que preservar}
    \small São o \highlight{baseline do varejo} e o denominador de todas as
    comparações já reportadas no Capítulo 5. Mudar sua fórmula invalidaria
    a leitura de qualquer resultado anterior que as usa como referência.
  \end{block}

  \note{
    Estas três são as únicas políticas do portfólio original de 12 que
    permanecem byte-a-byte equivalentes em lógica -- só a chamada ao
    simulador por baixo delas mudou (do InventoryEnv escalar, que
    continua existindo e sendo usado por elas, sem alteração). A decisão
    de mantê-las intocadas foi deliberada: são o piso de comparação usado
    em todo o Capítulo 5 da proposta, e qualquer mudança nelas quebraria a
    comparabilidade com os números já reportados.
  }
\end{frame}

% Slide — Novas: SOTA clássicas
\begin{frame}{\novo: três políticas clássicas em versão estado da arte}
  \footnotesize Complementam (não substituem) EOQ/$(s,S)$/Jornaleiro -- mesmo
  slot conceitual, referência publicada mais recente.

  \renewcommand{\arraystretch}{0.95}\scriptsize
  \begin{tabular}{@{}p{0.20\linewidth}p{0.52\linewidth}p{0.22\linewidth}@{}}
    \toprule
    \textbf{Política} & \textbf{Ideia central} & \textbf{Referência} \\
    \midrule
    PIL (\textit{Projected Inv.\ Level}) &
      $Q_t=\max(0, S-\mathrm{PIL}_t)$, estoque projetado $L$ ciclos à
      frente por reamostragem empírica, $\max(0,\cdot)$ em \textbf{cada}
      ciclo (evita viés sob \textit{lost sales}) &
      van Jaarsveld \& Arts, \textit{OR} 72(5), 2024 \\
    \addlinespace
    Capped base-stock &
      $Q_t=\min(\text{cap}, \max(0, S-\mathrm{IP}_t))$ -- híbrido
      base-stock/pedido constante; piso honesto p/ comparar RL &
      Xin, \textit{OR} 69(1), 2021 \\
    \addlinespace
    Big Data Newsvendor &
      Regressão quantílica, $\tau=c_s/(c_s{+}h)$, 8 características de
      intermitência (sem suposição de distribuição) &
      Ban \& Rudin, \textit{OR} 67(1), 2019 \\
    \bottomrule
  \end{tabular}

  \note{
    Estas três não existiam em julho -- é portfólio novo, não correção. O
    critério de escolha foi: para cada slot conceitual (banda de reposição,
    base-stock, jornaleiro), qual é a referência mais recente publicada em
    veículo relevante. O PIL corrige um viés sutil de projeção sob venda
    perdida: projetar linearmente sem truncar em zero a cada ciclo
    intermediário superestima a recuperação do estoque, o que é
    especialmente grave em séries Lumpy onde o estoque zera com
    frequência. O capped base-stock entra porque a literatura de RL para
    inventário usa exatamente essa política como piso -- incluí-la torna
    o benchmark honesto, evitando comparar agentes de RL só contra
    heurísticas fracas. O Big Data Newsvendor troca a suposição de
    distribuição normal por minimização de risco empírico via regressão
    quantílica, usando características de intermitência (fração de zeros
    recentes, tempo desde a última venda) que carregam mais sinal do que
    a média em regime Lumpy.
  }
\end{frame}

% Slide — Novas: Zabraoui
\begin{frame}{\novo: três heurísticas do artigo-base (Zabraoui et al., 2025)}
  \footnotesize Seção 3.5 do artigo-base define quatro políticas de
  referência; apenas $(s,S)$ já existia no portfólio. As outras três entram
  aqui.

  \renewcommand{\arraystretch}{0.95}\scriptsize
  \begin{tabular}{@{}p{0.24\linewidth}p{0.66\linewidth}@{}}
    \toprule
    \textbf{Política} & \textbf{Regra adotada} \\
    \midrule
    Min-Max &
      Pede até $M$ quando posição $\leq m$;
      $m=\mu L$ (lead time $L$), $M=m + \mu\cdot$cobertura \\
    Fixed Interval Repl. &
      Revisão a cada $R$ ciclos; repõe até
      $S=\mu(R{+}L)+z\sigma\sqrt{R{+}L}$ \\
    Vendor-Responsive &
      $s_t=\mu L + z_t\sigma\sqrt{L}$,
      $z_t=z_0(1+\kappa(\alpha_{\text{alvo}}-\mathrm{SL}_t))$ -- realimentação
      pelo serviço realizado \\
    \bottomrule
  \end{tabular}
  \vspace{0.3em}

  \begin{alertblock}{Ressalva de fidelidade, registrada no código}
    \scriptsize Min-Max e Fixed Interval: descritas no artigo em
    \textbf{uma frase cada, sem fórmula}. Vendor-Responsive é a mais
    subespecificada (``\textit{adapts order timing based on supplier lead
    times and historical service levels}''). Fórmulas acima são leitura
    razoável, \highlight{não transcrição literal} -- dito explicitamente
    nos comentários do código.
  \end{alertblock}

  \note{
    Este slide existe para ser transparente sobre o grau de fidelidade de
    cada implementação. O artigo do Zabraoui é a referência que este
    trabalho adota para o que não tem estado da arte próprio -- mas o
    artigo descreve essas três políticas em uma frase cada, sem equação.
    Documentei no próprio código, com essa mesma ressalva, qual foi a
    leitura escolhida: Min-Max e Fixed Interval seguem a definição
    clássica de manual de operações; Vendor-Responsive foi a mais livre de
    interpretar, porque "adapta o timing do pedido com base no lead time
    do fornecedor e no nível de serviço histórico" não tem parâmetro
    algum -- montei um laço de realimentação em que o z-score de segurança
    sobe quando o serviço realizado fica abaixo da meta. É importante que
    a banca saiba que isso é uma escolha nossa razoável, não uma
    reprodução certificada do artigo.
  }
\end{frame}


\section{Meta-heurísticas}
% Slide — GA: antes (julho)
\begin{frame}{Algoritmo Genético -- \antes}
  \footnotesize Biblioteca \texttt{DEAP}. Representação real $\theta=(\mathrm{ROP}, Q, \mathrm{SS})$.

  \renewcommand{\arraystretch}{1.05}\footnotesize
  \begin{tabular}{@{}p{0.28\linewidth}p{0.66\linewidth}@{}}
    \toprule
    \textbf{Operador} & \textbf{Implementação} \\
    \midrule
    População & 100 indivíduos, inicializados uniformemente nos limites de busca \\
    Gerações & 50 \\
    Seleção & Torneio, tamanho 3 (\texttt{tools.selTournament}) \\
    Cruzamento & \textit{Blend crossover} $\alpha=0{,}5$ (\texttt{tools.cxBlend}), prob. 0,8 \\
    Mutação & Gaussiana, $\sigma=$ maior amplitude do espaço de busca $/10$, prob. 0,05 por gene (\texttt{tools.mutGaussian}) \\
    Elitismo & \textbf{Nenhum} -- \texttt{HallOfFame(1)} só registra o melhor, não o protege do crossover \\
    Avaliação & 1 simulação completa de $T$ ciclos por indivíduo, \textbf{em laço Python} \\
    Aptidão & Soma ponderada $\lambda_1\,\mathrm{NS}-\lambda_2\,\mathrm{CTI}$, pesos fixos \\
    \bottomrule
  \end{tabular}

  \note{
    Este é o inventário completo do GA de julho, operador por operador,
    porque a pergunta natural da banca é "o que exatamente mudou". A
    população, o número de gerações, o tipo de cruzamento e o tipo de
    mutação continuam os mesmos na versão atual -- não houve mudança de
    metaheurística, só de execução e de aptidão. O ponto crítico aqui é a
    ausência de elitismo: sem proteção, o melhor indivíduo encontrado numa
    geração pode ser destruído pelo cruzamento blend na geração seguinte,
    o que explica parte da instabilidade observada entre séries nos
    resultados do Experimento 1. E cada geração custava 100 simulações
    sequenciais em Python puro.
  }
\end{frame}

% Slide — GA: depois (atual)
\begin{frame}{Algoritmo Genético -- \depois}
  \footnotesize PyTorch (\texttt{TorchGA}), sem DEAP. Mesma representação $\theta=(\mathrm{ROP}, Q, \mathrm{SS})$.

  \renewcommand{\arraystretch}{0.95}\scriptsize
  \begin{tabular}{@{}p{0.24\linewidth}p{0.72\linewidth}@{}}
    \toprule
    \textbf{Operador} & \textbf{Implementação} \\
    \midrule
    População & 100 (\textbf{inalterado}) \\
    Gerações & 50 na Bahia/"bot"; \textbf{10} no M5 (orçamento de tempo) \\
    Seleção & Torneio, tamanho 3 (\textbf{inalterado}), vetorizada:
      \texttt{torch.randint} sorteia $100\times3$ candidatos de uma vez,
      \texttt{argmin} por linha decide os vencedores \\
    Cruzamento & \highlight{Corrigido}: \textit{uniform crossover} -- cada
      gene herda do pai A ou B com $p=0{,}5$ (era \textit{blend}/BLX-$\alpha$,
      herdado do DEAP de julho) \\
    Mutação & \highlight{Corrigido}: gaussiana \textbf{adaptativa} -- taxa
      decai de 0,05 (1ª geração) até $0{,}05\times0{,}2$ (última), era taxa
      fixa \\
    Elitismo & \highlight{Novo}: melhor indivíduo histórico substitui o
      pior da nova geração -- nunca é perdido \\
    Avaliação & \highlight{Toda a população em 1 simulação vetorizada}
      (\texttt{BatchInventoryEnv}, largura 100) por geração \\
    Aptidão & Restrita (Eq.~4.2) \highlight{+ termos de risco} TR/BE \\
    \bottomrule
  \end{tabular}
  \vspace{0.3em}

  \begin{alertblock}{Por que o cruzamento e a mutação mudaram de novo}
    \scriptsize Auditoria (18/08) achou que o comentário do código citava
    "Zabraoui Tab.5" para o GA inteiro, mas só população/crossover\_prob/
    mutation\_prob batiam -- o artigo (Seç.\ 3.4/4.4) descreve
    \textit{uniform crossover} e \textit{adaptive mutation}, nunca
    implementados. Corrigido nesta sessão.
  \end{alertblock}

  \note{
    Este slide mudou duas vezes na mesma sessão de trabalho. A primeira
    versão comparava seleção/cruzamento/mutação como operadores idênticos
    ao slide anterior -- e isso era verdade em relação a julho, mas
    escondia um problema mais profundo: o comentário do código dizia
    "hiperparâmetros conforme Zabraoui Tabela 5" de um jeito que sugeria
    que o GA inteiro estava alinhado ao artigo, quando só os NÚMEROS
    (população, probabilidade de cruzamento, taxa inicial de mutação)
    batiam -- o MECANISMO dos operadores (blend crossover, mutação de
    taxa fixa) nunca foi trocado para o que o artigo realmente descreve
    (uniform crossover, mutação adaptativa). Auditado e corrigido nesta
    sessão: cruzamento uniforme (cada gene herda de um dos pais com
    probabilidade 0,5, não uma combinação linear contínua) e mutação com
    taxa decrescente ao longo das gerações (mais exploração no início,
    mais refinamento no fim -- leitura nossa de "adaptive mutation", o
    artigo não dá fórmula, documentada como tal no código). O elitismo
    continua sendo a correção mais simples e mais importante: sem ele, o
    cruzamento pode destruir a melhor solução encontrada geração após
    geração.
  }
\end{frame}

% Slide — SA antes/depois
\begin{frame}{Simulated Annealing -- \antes\ vs. \depois}
  \begin{columns}[T,onlytextwidth]
    \column{0.48\textwidth}
      \antes\\[0.3em]
      \small \texttt{scipy.optimize.dual\_annealing}
      \begin{itemize}\footnotesize
        \item Cronograma de temperatura e critério de aceitação
          \textbf{internos à biblioteca} -- caixa-preta com parâmetros
          próprios, não os documentados na proposta.
        \item Refinamento local por Nelder-Mead a cada iteração.
        \item Uma cadeia, avaliação sequencial.
      \end{itemize}
    \column{0.48\textwidth}
      \depois\\[0.3em]
      \small \texttt{TorchSA}: cronograma explícito
      \begin{itemize}\footnotesize
        \item $T_0=1000$, $\alpha_{\text{resfr.}}=0{,}95$ (Seção 4.4 da
          proposta), $T_{k+1}=\alpha\, T_k$.
        \item Critério de Metropolis auditável: aceita se $\Delta C<0$, ou
          com probabilidade $e^{-\Delta C/T_k}$.
        \item \highlight{32 cadeias independentes} (\texttt{n\_chains=32})
          em paralelo, mesmo cronograma, saída = melhor entre elas.
      \end{itemize}
  \end{columns}
  \vspace{0.8em}

  \begin{block}{Por que isso é uma reescrita, não só uma correção}
    \small SA é inerentemente sequencial. Vetorizar sem descaracterizar o
    método significa rodar várias cadeias \textbf{independentes} com o
    mesmo cronograma, em vez de fingir paralelismo dentro de uma cadeia.
    O ganho de robustez vem de reportar o melhor entre 32 buscas, não de
    uma cadeia mais rápida.
  \end{block}

  \note{
    O scipy.dual\_annealing é uma implementação respeitável, mas é uma
    caixa-preta: o cronograma de temperatura e os parâmetros de aceitação
    internos não são os mesmos documentados na Seção 4.4 da proposta, o
    que dificultava justificar por que aquele T0 e aquele alpha específicos.
    A versão atual implementa o critério de Metropolis manualmente, o que
    o torna auditável linha a linha, e explora paralelismo de um jeito
    metodologicamente honesto: como SA clássico é sequencial por natureza,
    a vetorização aqui significa rodar 32 cadeias independentes com o
    MESMO cronograma de temperatura, não inventar um SA "paralelo" que não
    existe na literatura -- a saída é simplesmente a melhor solução entre
    as 32.
  }
\end{frame}

% Slide — PSO antes/depois
\begin{frame}{Particle Swarm Optimization -- \antes\ vs. \depois}
  \begin{columns}[T,onlytextwidth]
    \column{0.48\textwidth}
      \antes\\[0.3em]
      \small NumPy puro, laço Python por partícula
      \begin{itemize}\footnotesize
        \item 40 partículas $\times$ 80 iterações.
        \item $w=0{,}7$, $c_1=c_2=1{,}5$, \textbf{sem} fator de constrição.
        \item Velocidade sem limite explícito -- pode divergir.
        \item Avaliação: laço \texttt{for i in range(n\_particles)}, uma
          simulação por partícula.
      \end{itemize}
    \column{0.48\textwidth}
      \depois\\[0.3em]
      \small \texttt{TorchPSO}, avaliação em lote
      \begin{itemize}\footnotesize
        \item 40 $\times$ 80 (\textbf{inalterado}).
        \item Mesmos $w$, $c_1$, $c_2$ \highlight{+ fator de constrição de
          Clerc \& Kennedy (2002)}:
          $\chi = \dfrac{2}{|2-\phi-\sqrt{\phi^2-4\phi}|}$, $\phi=c_1{+}c_2$.
        \item $v \leftarrow \chi\,[w v + c_1 r_1(p_{\text{best}}{-}x) +
          c_2 r_2(g_{\text{best}}{-}x)]$ -- converge sem \textit{clamp}
          arbitrário de velocidade.
      \end{itemize}
  \end{columns}
  \vspace{0.8em}

  \begin{alertblock}{Sintoma observado que motivou a correção}
    \small Nos resultados anteriores, $Q$ frequentemente ``colava'' no
    limite superior do espaço de busca -- assinatura clássica de
    velocidade divergente em PSO sem constrição.
  \end{alertblock}

  \note{
    O PSO antigo usava a combinação w=0,7 e c1=c2=1,5 sem nenhum mecanismo
    de controle de velocidade -- combinação em que, pela teoria de
    convergência de PSO, a velocidade PODE divergir. O sintoma prático era
    Q colando repetidamente no teto do espaço de busca, o que é
    exatamente o que se espera quando a velocidade cresce sem controle: a
    partícula "voa" para fora da região de interesse e é sempre truncada
    no limite. O fator de constrição de Clerc e Kennedy resolve isso
    matematicamente, sem precisar de um clamp arbitrário de velocidade que
    seria mais um parâmetro para justificar.
  }
\end{frame}

% Slide — DE antes/depois
\begin{frame}{Differential Evolution -- \antes\ vs. \depois}
  \begin{columns}[T,onlytextwidth]
    \column{0.48\textwidth}
      \antes\\[0.3em]
      \small \texttt{scipy.optimize.differential\_evolution}
      \begin{itemize}\footnotesize
        \item \texttt{popsize=15} (multiplicador interno do scipy $\to$
          população efetiva maior, não documentada explicitamente).
        \item Estratégia padrão do scipy (\texttt{best1bin}), mutação em
          faixa $(0{,}5,\,1{,}0)$.
        \item \texttt{polish=True}: refinamento local por L-BFGS-B ao final
          -- passo extra fora da formulação DE canônica.
      \end{itemize}
    \column{0.48\textwidth}
      \depois\\[0.3em]
      \small \texttt{TorchDE}: DE/rand/1/bin explícito (Storn \& Price, 1997)
      \begin{itemize}\footnotesize
        \item População $=15\times3=45$ (\textbf{tamanho equivalente},
          agora explícito no código, sem multiplicador implícito).
        \item Mutação: $v_i = x_{r_1} + F(x_{r_2}{-}x_{r_3})$, $F=0{,}8$.
        \item Cruzamento binomial com $CR=0{,}9$, garantindo $\geq1$
          dimensão herdada do mutante (\texttt{jrand}).
        \item Seleção \textbf{gulosa} elemento a elemento -- \textbf{sem}
          polimento local: fiel à formulação DE/rand/1/bin canônica.
      \end{itemize}
  \end{columns}
  \vspace{0.6em}

  \small Toda a geração avaliada em uma chamada; a seleção gulosa é um
  único \texttt{torch.where} -- não há laço Python sobre indivíduos.

  \note{
    O DE antigo delegava a uma implementação de caixa-preta do scipy, cuja
    estratégia padrão (best1bin) e cujo passo de polimento local por
    L-BFGS-B ao final não são parte da formulação DE/rand/1/bin canônica
    de Storn e Price -- eram comportamento padrão do scipy, não uma
    escolha metodológica documentada. A versão atual implementa a
    formulação canônica explicitamente: mutação diferencial rand/1,
    cruzamento binomial garantindo que pelo menos uma dimensão venha do
    mutante, e seleção gulosa pura, sem qualquer refinamento local
    adicional. Isso torna o método auditável e consistente com a citação
    de Storn \& Price (1997) usada na dissertação.
  }
\end{frame}


\section{A função de aptidão}
% Slide — Problema de escala na aptidão antiga
\begin{frame}{A função de aptidão tem sua própria evolução}
  \textbf{\textcolor{darkblue}{Três versões, em ordem cronológica -- não é uma correção pontual, é uma linha evolutiva.}}
  \vspace{0.5em}

  \begin{enumerate}
    \item \textbf{\antes} -- soma ponderada (Eq.~3 do artigo-base):
      $\mathrm{Fitness}=\lambda_1\,\mathrm{NS}-\lambda_2\,\mathrm{CTI}$,
      $(\lambda_1,\lambda_2)=(1{,}0;\ 0{,}0001)$ fixos.
    \item \textbf{Intermediário (ago/2026, reimplementação)} -- formulação
      restrita (Eq.~4.2 da proposta): $\min \mathrm{CTI}$ s.a.
      $\mathrm{NS}\geq\alpha_{\min}$, penalidade relativa ao próprio CTI.
    \item \textbf{\novo\ (esta sessão)} -- qualquer um dos dois modos
      \highlight{$+$ termos de risco} baseados em TR e BE (próximos slides).
  \end{enumerate}
  \vspace{0.6em}

  \begin{alertblock}{Por que a versão 1 é problemática}
    \small O CTI varia por \textbf{ordens de grandeza} entre séries (razão
    de $124\times$ entre a maior e a menor loja no estudo piloto). Pesos
    fixos fazem a MESMA configuração privilegiar serviço em séries de
    baixo volume e custo em séries de alto volume -- contaminando a
    comparação entre políticas.
  \end{alertblock}

  \note{
    É importante deixar claro que a função de aptidão passou por uma
    evolução em três estágios, não uma correção única. A primeira versão é
    literalmente a Equação 3 do artigo-base do Zabraoui -- reproduzida
    fielmente, com os mesmos pesos. O problema é que essa soma ponderada
    funciona bem no contexto do artigo (itens de alto giro, M5, escala
    homogênea), mas no recorte brasileiro, onde o CTI varia 124 vezes
    entre a maior e a menor loja, o mesmo par de pesos produz
    comportamentos completamente diferentes dependendo do volume da
    série. A segunda versão resolve isso tornando a penalidade relativa ao
    próprio custo. A terceira, mais recente, adiciona dois termos de risco
    que nem a formulação restrita nem a soma ponderada capturavam.
  }
\end{frame}

% Slide — Formulação restrita
\begin{frame}{Estágio 2: formulação restrita (Eq.~4.2)}
  \footnotesize
  \begin{columns}[T,onlytextwidth]
    \column{0.5\textwidth}
      \textbf{Objetivo, a minimizar:}
      \[
        \mathrm{custo} = \mathrm{CTI} + \delta \cdot p \cdot
        \max(\mathrm{CTI}, 1)
      \]
      \[
        \delta = \max(0,\ \alpha_{\min} - \mathrm{NS})
      \]
      em que $\alpha_{\min}=0{,}70$ e $p=10$ (peso de penalidade).
    \column{0.5\textwidth}
      \textbf{Por que isso resolve o problema de escala}
      \begin{itemize}
        \item A penalidade é \highlight{proporcional ao próprio CTI}
          (via $\max(\mathrm{CTI},1)$), não uma constante fixa.
        \item Candidatos viáveis ($\mathrm{NS}\geq\alpha_{\min}$) são
          ordenados por $-\mathrm{CTI}$ puro -- sem distorção de escala.
        \item Mesma fórmula usada nas versões escalar
          (\texttt{policies.\_eval\_static}) e vetorizada
          (\texttt{constrained\_cost}) -- \textbf{idênticas por construção}.
      \end{itemize}
  \end{columns}
  \vspace{0.8em}

  \small \textbf{Retrocompatibilidade:} \texttt{fitness\_mode="weighted"}
  preserva o comportamento antigo, para reproduzir os resultados já
  reportados no Capítulo 5 se necessário.

  \note{
    A formulação restrita segue exatamente a Equação 4.2 da proposta:
    minimizar CTI sujeito a NS maior ou igual a 0,70. Em vez de descartar
    candidatos inviáveis (o que perderia informação de gradiente para a
    busca), a restrição entra como penalidade proporcional ao déficit de
    serviço, escalada pelo próprio CTI da trajetória. Isso é o que resolve
    o problema de escala: como a penalidade é relativa, não absoluta, o
    mesmo critério funciona igualmente bem numa loja de CTI=50 e numa de
    CTI=6000. O modo antigo continua disponível por retrocompatibilidade,
    caso seja necessário reproduzir exatamente os números já publicados no
    capítulo de resultados.
  }
\end{frame}

% Slide — Termos de risco TR/BE
\begin{frame}{\novo: termos de risco (TR, BE) somados à aptidão}
  \footnotesize
  \textbf{\textcolor{darkblue}{CTI só registra o custo de ruptura JÁ REALIZADO na trajetória simulada. TR e BE capturam fragilidade que o CTI sozinho não pune.}}
  \vspace{0.4em}

  \[
    \mathrm{fitness}\ \mathrel{+}=\ \underbrace{\lambda_{\mathrm{TR}}\cdot \mathrm{TR}\cdot(1-\mathrm{NS})\cdot \mathrm{CTI}_{\mathrm{ref}}}_{\text{risco de ruptura}}
    \ +\ \underbrace{\lambda_{\mathrm{BE}}\cdot \max(\mathrm{BE}-1,\,0)\cdot \mathrm{CTI}_{\mathrm{ref}}}_{\text{compra excessiva}}
  \]

  \begin{columns}[T,onlytextwidth]
    \column{0.5\textwidth}
      \textbf{TR $\times$ (1$-$NS)}
      \begin{itemize}
        \item TR = fração de ciclos em que \highlight{falta produto pro
          cliente} (probabilidade de ruptura).
        \item Modulado por $(1{-}\mathrm{NS})$: risco composto cresce
          quando serviço \textbf{já ruim} E ruptura \textbf{frequente}
          coincidem.
      \end{itemize}
    \column{0.5\textwidth}
      \textbf{$\max(\mathrm{BE}-1,0)$}
      \begin{itemize}
        \item $\mathrm{BE}=\mathrm{Var}(Q)/\mathrm{Var}(D)$.
        \item $\mathrm{BE}\leq1$: pedido tão ou mais suave que a demanda
          -- \textbf{não punido}.
        \item $\mathrm{BE}>1$: amplificação -- só o \highlight{excesso}
          acima de 1 é punido (``compra excessiva'').
      \end{itemize}
  \end{columns}
  \vspace{0.6em}
  \small $\lambda_{\mathrm{TR}}=\lambda_{\mathrm{BE}}=1{,}0$ por padrão;
  $\mathrm{CTI}_{\mathrm{ref}}=\max(\mathrm{CTI},1)$, mesmo truque de escala
  da penalidade de NS. Soma-se a QUALQUER \texttt{fitness\_mode}
  (\texttt{constrained} ou \texttt{zabraoui}).

  \note{
    Esta é a extensão metodológica mais recente, feita nesta mesma sessão
    de acompanhamento. A motivação veio de uma auditoria de onde TR e BE
    eram usados no pipeline: eram calculadas e reportadas como KPIs, mas
    NUNCA entravam na busca de nenhuma meta-heurística -- o CTI sozinho
    define a busca, e o CTI só reflete o custo de ruptura que de fato
    aconteceu na trajetória simulada, não o RISCO de que aconteça de novo.
    TR entra como uma probabilidade de gerar custo futuro, ponderada por
    NS porque o pior cenário é quando as duas coisas coincidem: serviço já
    ruim e ruptura frequente. BE entra como punição especificamente pela
    amplificação do pedido sobre a variabilidade da própria demanda -- o
    "efeito chicote" -- mas só acima de 1, porque BE menor ou igual a 1 é
    comportamento saudável, não deve ser punido. A formulação foi validada
    contra a literatura: a definição de BE como razão de variâncias com
    1,0 como zero natural vem da literatura padrão de bullwhip effect; a
    penalidade por probabilidade de ruptura em otimização com restrição de
    nível de serviço tem precedente publicado em pesquisa operacional; e
    há precedente de GA com aptidão penalizada especificamente para
    mitigar bullwhip effect na literatura de cadeia de suprimentos.
  }
\end{frame}


\section{Aprendizado por reforço}
% Slide — DQN antes/depois
\begin{frame}{DQN -- \antes\ vs. \depois}
  \begin{columns}[T,onlytextwidth]
    \column{0.48\textwidth}
      \antes\\[0.3em]
      \small Rede \texttt{\_NN}: NumPy escrito à mão
      \begin{itemize}\footnotesize
        \item MLP simples, \textit{forward}/\textit{backward} manuais,
          \texttt{ReLU}, sem otimizador adaptativo (SGD puro).
        \item DQN padrão: mesma rede escolhe e avalia a ação do próximo
          estado -- \textbf{superestimação} de valor conhecida na
          literatura.
        \item Perda: erro quadrático médio.
        \item Sem \textit{clipping} de gradiente.
      \end{itemize}
    \column{0.48\textwidth}
      \depois\\[0.3em]
      \small \texttt{DoubleDQNAgent} (PyTorch)
      \begin{itemize}\footnotesize
        \item Arquitetura \highlight{\textit{dueling}} (Wang et al., 2016):
          $Q(s,a)=V(s)+[A(s,a)-\overline{A(s,\cdot)}]$.
        \item \highlight{Double DQN} (van Hasselt et al., 2016): rede
          \textit{online} escolhe a ação, rede \textit{alvo} avalia --
          $y=r+\gamma\,Q_{\text{alvo}}(s',\arg\max_a Q_{\text{online}}(s',a))$.
        \item Perda de \textbf{Huber} (\texttt{smooth\_l1\_loss}), Adam,
          \textit{grad clipping} em 10.
      \end{itemize}
  \end{columns}
  \vspace{0.6em}

  \begin{block}{Por que Double DQN dueling especificamente}
    \small Em horizontes curtos ($T=38$ ciclos), a superestimação do DQN
    clássico é especialmente danosa: o agente aprende a pedir demais
    porque valoriza excessivamente estados de estoque alto. A separação
    $V(s)/A(s,a)$ estabiliza o treino em estados onde a ação pouco altera
    o retorno -- frequente quando o estoque já está muito acima do ponto
    de pedido.
  \end{block}

  \note{
    O que era rotulado "DQN" em julho era, tecnicamente, DQN vanilla com
    uma rede neural implementada do zero em NumPy -- funcional, mas sem
    nenhuma das melhorias que a literatura de RL desenvolveu desde 2015
    especificamente para o problema de superestimação de valor. A versão
    atual não é apenas uma reescrita em PyTorch: é a adoção de Double DQN
    com arquitetura dueling, que são, juntos, considerados a versão de
    estado da arte razoável para este slot no portfólio -- van Hasselt et
    al. 2016 para o Double DQN, Wang et al. 2016 para dueling. A perda de
    Huber, em vez de erro quadrático, também é mais robusta a
    recompensas de escala variável entre séries.
  }
\end{frame}

% Slide — PPO: o bug em detalhe
\begin{frame}{PPO -- o defeito, em detalhe matemático}
  \footnotesize
  \textbf{Objetivo correto do PPO (Schulman et al., 2017):}
  \[
    L(\theta) = -\mathbb{E}\Big[\min\big(r_t A_t,\ \mathrm{clip}(r_t, 1{-}\epsilon, 1{+}\epsilon)\,A_t\big)\Big],
    \qquad r_t = \frac{\pi_{\text{novo}}(a_t|s_t)}{\pi_{\text{antigo}}(a_t|s_t)}
  \]
  \vspace{0.3em}

  \begin{columns}[T,onlytextwidth]
    \column{0.48\textwidth}
      \antes\ (o defeito)
      \begin{itemize}\footnotesize
        \item \texttt{tgt\_a[i, A[i]] += adv[i] * min(rat[i], cl[i])}
        \item Soma a vantagem \textbf{diretamente aos logits-alvo} e treina
          o ator por \textbf{erro quadrático} contra esse alvo.
        \item Isso \highlight{NÃO é o gradiente} de $L(\theta)$ acima --
          é uma heurística sem garantia de melhoria monotônica da política.
      \end{itemize}
    \column{0.48\textwidth}
      \depois\ (a correção)
      \begin{itemize}\footnotesize
        \item Gradiente obtido por \textbf{autograd} diretamente sobre
          $r_t A_t$ e o termo recortado, como na formulação original.
        \item \texttt{policy\_loss = -torch.min(s1, s2).mean()}
        \item \highlight{GAE$(\lambda{=}0{,}95)$} para o estimador de
          vantagem, normalização de vantagem, recorte também da função
          valor, bônus de entropia.
      \end{itemize}
  \end{columns}
  \vspace{0.6em}

  \begin{alertblock}{Evidência do impacto}
    \small Experimento 1: PPO com FP $=0{,}98$ (pede em quase todo ciclo)
    -- comportamento consistente com um ator que nunca convergiu para uma
    política seletiva, porque nunca foi de fato otimizado pelo objetivo
    do PPO.
  \end{alertblock}

  \note{
    Este é o defeito mais grave encontrado na auditoria, porque não é uma
    questão de estado da arte ou de escolha de hiperparâmetro -- é
    matematicamente incorreto. Somar a vantagem aos logits-alvo e treinar
    por MSE é uma heurística que parece plausível ("empurra a
    probabilidade da ação boa para cima"), mas não corresponde ao
    gradiente de nenhuma função objetivo bem definida, muito menos à
    surrogate function recortada do PPO, que é o que garante que a
    atualização não afaste demais a nova política da antiga. Sem essa
    garantia, o treino pode divergir de formas imprevisíveis -- e o
    sintoma observado, frequência de pedido de 0,98, é exatamente o tipo
    de degeneração esperada: o agente não aprendeu a diferenciar quando
    pedir e quando não pedir, então pede quase sempre. A correção usa
    autograd do PyTorch sobre a expressão matemática correta, o que
    elimina a possibilidade desse tipo de erro de derivação manual.
  }
\end{frame}

% Slide — SARSA antes/depois
\begin{frame}{SARSA -- \antes\ vs. \depois}
  \begin{columns}[T,onlytextwidth]
    \column{0.48\textwidth}
      \antes\\[0.3em]
      \small SARSA tabular clássico
      \begin{itemize}\footnotesize
        \item Tabela $Q$ de $5\times10$ (5 estados de estoque $\times$ 10
          ações).
        \item Discretização \textbf{1D}: só nível de estoque normalizado.
        \item Alvo TD amostra a próxima ação $a'$ realmente tomada
          ($\epsilon$-greedy) -- variância de amostragem.
        \item $\alpha=0{,}1$, $\gamma=0{,}99$, $\epsilon=0{,}1$.
      \end{itemize}
    \column{0.48\textwidth}
      \depois\\[0.3em]
      \small \texttt{ExpectedSARSAAgent}
      \begin{itemize}\footnotesize
        \item Tabela $Q$ de $(5\times3)\times10$ -- discretização
          \highlight{\textbf{2D}}: estoque $\times$ pedidos em trânsito
          (\texttt{n\_pipeline\_states=3}).
        \item Alvo = \highlight{esperança} sob a política $\epsilon$-greedy:
          $y=r+\gamma\sum_a \pi(a|s')\,Q(s',a)$, não uma amostra.
        \item Mesmos $\alpha$, $\gamma$, $\epsilon$ (\textbf{inalterados}).
      \end{itemize}
  \end{columns}
  \vspace{0.8em}

  \begin{block}{Sem estado da arte publicado para este slot}
    \small Diferente de DQN e PPO, não existe estado da arte publicado em
    veículo relevante para SARSA aplicado a inventário. A escolha foi
    Expected SARSA (van Seijen et al., ADPRL 2009) -- reduz a variância do
    SARSA amostrado mantendo o caráter \textit{on-policy}, permanecendo
    deliberadamente como \textbf{baseline tabular simples} do portfólio.
  \end{block}

  \note{
    O SARSA é o único slot de RL em que não há uma referência mais recente
    óbvia para adotar -- por isso a mudança aqui é mais modesta que em DQN
    e PPO. Duas melhorias, ambas bem estabelecidas: Expected SARSA
    substitui a amostra da próxima ação por sua esperança sob a política
    corrente, o que elimina uma fonte de variância sem abandonar o
    caráter on-policy que distingue SARSA de Q-learning; e a
    discretização do estado ganhou uma segunda dimensão -- pedidos em
    trânsito -- porque a versão anterior, que só olhava para o nível de
    estoque, não conseguia diferenciar um estado de estoque baixo com
    pedido a caminho de um estado de estoque baixo sem nada em trânsito,
    que exigem decisões diferentes.
  }
\end{frame}


\section{Híbridas GA-RL}
% Slide — Híbridas GA-RL antes/depois
\begin{frame}{GA-DQN / GA-PPO -- \antes\ vs. \depois}
  \small GA decide \textbf{QUANDO} pedir (limiar $\mathrm{ROP}{+}\mathrm{SS}$); o agente de RL decide \textbf{QUANTO} -- inalterado nas duas versões.
  \vspace{0.4em}

  \begin{columns}[T,onlytextwidth]
    \column{0.48\textwidth}
      \antes
      \begin{itemize}\footnotesize
        \item \texttt{qty = max(rl.act(s)[1], self.Q * 0.5)}
        \item Piso de \textbf{50\% do $Q$ do GA} fixo, sem parametrização.
        \item RL de baixo nível ainda sofria do defeito do PPO
          (Seção~6) quando o híbrido usava GA-PPO.
      \end{itemize}
    \column{0.48\textwidth}
      \depois
      \begin{itemize}\footnotesize
        \item Mesmo piso, agora \highlight{parametrizado}:
          \texttt{floor\_ratio} (padrão 0,5), configurável por
          \texttt{simulation.hybrid}.
        \item Herda \textbf{automaticamente} as correções de DQN e PPO
          (Seção~6) -- a classe \texttt{HybridGARL} recebe qualquer agente
          já treinado como componente.
        \item Herda também o GA vetorizado com elitismo (Seção~4) na fase
          (i) de inicialização.
      \end{itemize}
  \end{columns}
  \vspace{0.8em}

  \begin{alertblock}{Por que o piso de quantidade existe}
    \small Sem ele, o agente de RL pode aprender a emitir pedidos
    \textbf{simbólicos} (quantidade $\approx0$) só para satisfazer o
    gatilho de \textit{quando} pedir do GA -- comportamento degenerado
    observado no Experimento 1 com PPO sem inicialização adequada.
  \end{alertblock}

  \note{
    A arquitetura híbrida em si -- GA decide quando, RL decide quanto -- não
    mudou. O que mudou é que ela herda automaticamente todas as correções
    das duas peças que a compõem: o GA vetorizado com elitismo entra na
    fase de inicialização/pré-treino (prepopulate\_from\_ga para o DQN,
    warmstart\_from\_ga para o PPO), e o agente de RL de baixo nível agora é
    o Double DQN ou o PPO corrigido, não as versões antigas. O piso de
    quantidade (floor\_ratio) passou de uma constante hardcoded para um
    parâmetro configurável, o que facilita qualquer análise de
    sensibilidade futura sobre esse valor.
  }
\end{frame}


\section{Infraestrutura e hiperparâmetros}
% Slide — GPU routing
\begin{frame}{Sobre GPU: implementada, mas raramente usada}
  \keymessage{A vetorização é o que traz o ganho de desempenho -- não a GPU. E isso foi medido, não assumido.}
  \vspace{0.6em}

  \begin{itemize}
    \item \texttt{core/device.py}: roteamento CPU/CUDA/MPS
      \textbf{consciente do tamanho do lote} -- decide o acelerador só
      acima de um limiar medido, não por preferência fixa.
    \item Medição: MPS (Apple GPU) só compensa acima de
      \textbf{$\sim$150 mil trajetórias simultâneas}.
    \item Com população de 100 indivíduos (GA) ou 40 partículas (PSO), a
      \textbf{CPU é 35$\times$ mais rápida} -- overhead de transferência
      CPU$\leftrightarrow$GPU domina em lotes pequenos.
  \end{itemize}
  \vspace{0.8em}

  \begin{block}{Decisão de engenharia}
    \small Rotas de RL (\texttt{rl\_torch.py}) fixam CPU por padrão
    (redes de $\sim10^4$ parâmetros, lotes de 64 -- overhead de
    sincronização de GPU não compensa). Meta-heurísticas
    (\texttt{metaheuristics\_torch.py}) consultam o limiar medido a cada
    execução, por tamanho de população.
  \end{block}

  \note{
    Este slide existe para não deixar a impressão de que "PyTorch" implica
    "GPU". A infraestrutura de roteamento foi implementada e testada, mas
    a medição honesta mostrou que, na escala de população usada neste
    projeto -- dezenas a centenas de indivíduos por geração -- o overhead
    de transferência de dados entre CPU e GPU é maior que o ganho de
    paralelismo da GPU. O ganho real de desempenho, que tornou viável
    rodar sobre bases de dezenas de milhares de séries, veio inteiramente
    da vetorização (população inteira em uma chamada), não do hardware.
    Isso é consistente com a literatura de sistemas: GPUs compensam
    quando o paralelismo é massivo (aqui, só acima de 150 mil trajetórias
    simultâneas), não em lotes de 100.
  }
\end{frame}

% Slide — Hiperparâmetros alinhados ao Zabraoui (+ auditoria de fidelidade)
\begin{frame}{Hiperparâmetros: alinhados ao artigo-base -- e auditados}
  \footnotesize Onde \textbf{não existe} estado da arte publicado para
  meta-heurísticas em inventário, os hiperparâmetros foram alinhados às
  Tabelas 4 e 5 de Zabraoui et al. (2025). Uma auditoria de fidelidade
  (18/08) achou que parte disso ficava só no comentário do código, sem
  chegar no valor real.

  \renewcommand{\arraystretch}{0.95}\scriptsize
  \begin{tabular}{@{}p{0.24\linewidth}p{0.30\linewidth}p{0.38\linewidth}@{}}
    \toprule
    \textbf{Parâmetro} & \textbf{Citado como adotado} & \textbf{Valor real antes da auditoria} \\
    \midrule
    GA crossover\_prob & 0,8 (Tab.5) & 0,8 -- \textcolor{aftergreen}{já batia} \\
    GA mutation\_prob (inicial) & 0,05 (Tab.5) & 0,05 -- \textcolor{aftergreen}{já batia} \\
    GA crossover (mecanismo) & \textit{uniform} (Seç.3.4) & \textit{blend/BLX-$\alpha$} -- \textcolor{beforered}{divergia} \\
    GA mutação (mecanismo) & adaptativa (Seç.3.4) & taxa fixa -- \textcolor{beforered}{divergia} \\
    PPO entropy\_coef & 0,005 (Tab.4) & 0,005 -- \textcolor{aftergreen}{já batia} \\
    DQN epsilon & 0,2 fixo (Tab.4) & 1,0$\to$0,01 decrescente -- \textcolor{beforered}{divergia} \\
    Episódios DQN/PPO & $\geq$1000 (Seç.3.8) & 500 (50 no M5) -- \textcolor{beforered}{divergia} \\
    \bottomrule
  \end{tabular}
  \vspace{0.4em}

  \begin{alertblock}{Corrigido nesta sessão (não só a citação -- o código)}
    \scriptsize Os quatro itens que divergiam foram corrigidos no
    \texttt{simulation.yml} e no \texttt{TorchGA}: epsilon fixo, 1000
    episódios (100 no M5, mesmo corte de tempo de 10$\times$), cruzamento
    uniforme, mutação adaptativa. \highlight{Consequência}: as execuções de
    "bot"/M5 em andamento no momento da auditoria usavam os valores
    antigos e precisaram ser relançadas.
  \end{alertblock}

  \note{
    Este slide existe em duas camadas. A primeira -- por que esses valores
    específicos -- já estava no deck original: para meta-heurísticas em
    inventário não existe estado da arte publicado, então alinhar ao
    artigo-base é mais defensável que valores ad-hoc. A segunda camada,
    adicionada depois de uma auditoria pedida pelo usuário durante esta
    mesma sessão de trabalho, é mais desconfortável de admitir: nem tudo
    que o comentário do código citava como "adotado de Zabraoui" de fato
    tinha chegado ao valor configurado. DQN com epsilon decrescente em vez
    de fixo, 500 episódios em vez de 1000, cruzamento blend em vez de
    uniforme, mutação de taxa fixa em vez de adaptativa -- quatro casos em
    que a citação ficou no comentário e o config divergia silenciosamente.
    O usuário decidiu corrigir o código, não só a documentação, mesmo
    sabendo que isso invalidava as execuções de "bot" e M5 que estavam
    rodando havia quase duas horas no momento do pedido. É importante
    apresentar isso com transparência: não é vergonha ter encontrado o
    problema, seria vergonha não corrigir ou não relatar.
  }
\end{frame}


\section{Portfólio final}
% Slide — Portfólio final: 12 -> 18
\begin{frame}{Portfólio final: de 12 para 18 políticas}
  \centering
  \resizebox{\linewidth}{!}{%
  \tabstyle
  \begin{tabular}{@{}llll@{}}
    \toprule
    \textbf{Clássicas} & \textbf{SOTA / Zabraoui \novo} & \textbf{Busca \& RL} & \textbf{Híbridas} \\
    \midrule
    EOQ & PIL & GA & GA-DQN \\
    $(s,S)$ & CappedBaseStock & SA & GA-PPO \\
    Jornaleiro & BigDataNewsvendor & PSO & \\
     & Min-Max & DE & \\
     & Fixed Interval & DQN & \\
     & Vendor-Responsive & PPO & \\
     & & SARSA & \\
    \midrule
    \textbf{3} & \textbf{6 novas} & \textbf{7} & \textbf{2} \\
    \bottomrule
  \end{tabular}}
  \vspace{1em}

  \small \textbf{3 inalteradas} $\cdot$ \textbf{6 novas} $\cdot$
  \textbf{4 corrigidas} (GA, SA, PSO, DE por vetorização/aptidão; PPO por
  gradiente) $\cdot$ \textbf{5 reescritas} sem defeito conhecido (DQN,
  SARSA, GA-DQN, GA-PPO por herança das correções, simulador por
  vetorização).

  \note{
    O número final é 18: as 12 políticas de julho, menos zero (nenhuma foi
    removida), mais as 6 novas -- 3 do slot de estado da arte clássico
    (PIL, capped base-stock, big data newsvendor) e 3 do artigo-base
    Zabraoui (Min-Max, Fixed Interval, Vendor-Responsive). Cada célula
    deste portfólio tem agora uma justificativa documentada: ou é
    baseline preservado por comparabilidade, ou é a referência mais
    recente publicada para aquele slot conceitual, ou -- quando não existe
    referência publicada -- é a formulação canônica do método com a
    função objetivo corrigida.
  }
\end{frame}

% Slide — Próximos passos
\begin{frame}{Onde isso está sendo validado agora}
  \small As 18 políticas, com a função de aptidão em seu estágio mais
  recente (Eq.~4.2/Zabraoui $+$ TR/BE), estão sendo reexecutadas em
  \textbf{três bases}, com o mesmo critério de otimização em todas:
  \vspace{0.4em}

  \tabstyle\footnotesize
  \begin{tabular}{@{}lll@{}}
    \toprule
    \textbf{Base} & \textbf{Escopo} & \textbf{Status} \\
    \midrule
    Bahia (Experimento 2, oficial) & 145 séries & concluída \\
    ``bot'' (interna completa, 27 estados) & 4.869 séries & em andamento \\
    M5 (Walmart, benchmark externo, sem cortes) & 30.490 séries & em andamento \\
    \bottomrule
  \end{tabular}
  \vspace{1em}

  \begin{block}{Por que reexecutar a Bahia}
    \small A extensão TR/BE (estágio 3 da aptidão) foi aplicada
    \textbf{durante} a execução anterior da Bahia -- deixando etapas já
    concluídas (clássica, meta-heurística, híbrida) com a aptidão antiga.
    Reexecutada do zero para que as \textbf{três bases otimizem pelo mesmo
    critério}, condição necessária para qualquer comparação entre elas.
  \end{block}

  \note{
    Este slide fecha o círculo entre a implementação apresentada e o
    trabalho em andamento. Como a extensão de aptidão com TR e BE foi a
    mudança mais recente, aplicada no meio de uma execução de produção da
    Bahia, essa execução específica ficou com etapas inconsistentes entre
    si -- as que já tinham rodado usaram o critério antigo. A decisão foi
    reexecutar a Bahia do zero com o critério novo do início ao fim, para
    garantir que qualquer comparação futura entre Bahia, "bot" e M5 esteja
    otimizando exatamente o mesmo objetivo nas três. Isso não é
    apresentado como resultado -- é o estado do trabalho de validação no
    momento desta apresentação.
  }
\end{frame}


\end{document}
